\begin{proof}
    Assume that conditions (1) and (2) hold and define an assignment $F: P( \Sigma ) \to P( \Sigma ) / E$ such that $F_0( n ) = n + 1$ and $F( f : n \to n ) = \pi_E( \tau_n( f ) )$.
    Then $F$ is fully determined, and is therefore unique.
    It remains to be shown that $F$ is a functor.
    \begin{itemize}
    \item \textbf{Identities}.
          Let $n \in \mathbb{N}$.
          Since $\tau_n( 1_n ) \approx_E 1_{n+1}$ by condition (1), then the following holds.
          \begin{equation*}
              F( 1_n )
              =
              \pi_E( \tau_n( 1_n ) )
              =
              \pi_E( 1_{n+1} )
              =
              1_{n+1}
          \end{equation*}
          Since $n$ was arbitrary, then $F$ respects identities.
    \item \textbf{Composition}.
          Let $f: n \to n$ and $g: n \to n$.
          Since $\tau_n( g \circ f ) \approx_E \tau_n( g ) \circ \tau_n( f )$ by condition (2), then the following holds.
          \begin{equation*}
              F( g \circ f )
              =
              \pi_E( \tau_n( g \circ f ) )
              =
              \pi_E( \tau_n( g ) \circ \tau_n( f ) )
              =
              \pi_E( \tau_n( g ) ) \circ \pi_E( \tau_n( f ) )
              =
              F( g ) \circ F( f )
          \end{equation*}
          Since $f$ and $g$ were arbitrary, then $F$ respects composition.
    \end{itemize}
    Since $F$ respects identities and composition, then $F$ is a functor.
    Next, assume that condition (3) holds.
    It must be shown that there exists a unique functor $H: P( \Sigma ) / E \to P( \Sigma ) / E$ such that $H \circ \pi_E = F$.
    Let $( f, g ) \in E$.
    Since $\tau_E( f ) \approx_E \tau_E( g )$, then $F( f ) = \pi_E( \tau_n( f ) ) = \pi_E( \tau_n( g ) ) = F( g )$.
    Since $( f, g )$ was arbitrary, then there exists a unique functor $H: P( \Sigma ) / E \to P( \Sigma ) / E$ such that $H \circ \pi_E = F$.
    Finally, assume that conditions (4) through to (6) hold.
    It remains to be shown that $H$ is a control functor.
    \begin{enumerate}
    \item Let $f \in ( P( \Sigma ) / E )( n, n )$.
          Then there exists a $g \in P( \Sigma )( n, n )$ such that $\pi_E( g ) = f$.
          Then $H( f \boxtimes 1 ) = H( \pi_E( g ) \boxtimes 1 ) = H( \pi_E( g \boxtimes 1 ) ) = F( g \boxtimes 1)$.
          Since $\tau_{n+1}( g \boxtimes 1 ) \approx_E \tau_n( g ) \boxtimes 1$ by condition (4), then the following holds.
          \begin{equation*}
            H( f \boxtimes 1 )
            =
            F( g \boxtimes 1 )
            =
            \pi_E( \tau_n( g \boxtimes 1 ) )
            =
            \pi_E( \tau_n( g ) \boxtimes 1  )
            =
            \pi_E( \tau_n( g ) ) \boxtimes 1
            =
            F( g ) \boxtimes 1
            =
            H( f ) \boxtimes 1
          \end{equation*}
          Since $f$ was arbitrary, then $H( f \boxtimes 1 ) = H( f ) \boxtimes 1$ for each $f \in ( P( \Sigma ) / E )( n, n )$.
    \item Let $f \in ( P( \Sigma ) / E )( n, n )$.
          Then there exists a $g \in P( \Sigma )( n, n )$ such that $\pi_E( g ) = f$.
          Then the following equation holds.
          \begin{equation*}
              H( H( f ) )
              =
              H( H( \pi_E( g ) ) )
              =
              H( F( g ) )
              =
              H( \pi_E( \tau_n( g ) ) )
              =
              F( \tau_n( g ) )
              =
              \pi_E( \tau_{n+1}( \tau_n( g ) ) )
          \end{equation*}
          Since $\tau_{n+1}( \tau_n( g ) ) \approx_E ( \sigma \boxtimes 1_n ) \circ \tau_{n+1}( \tau_n( g ) ) \circ ( \sigma \boxtimes 1_n )$ by condition (5), then the following equation holds where $\sigma = \sigma \boxtimes 1_n$ and $\overline{\sigma} = \pi_E( \sigma )$.
          \begin{equation*}
              \pi_E( \tau_{n+1}( \tau_n( g ) ) )
              =
              \pi_E( \sigma \circ \tau_{n+1}( \tau_n( g ) ) \circ \sigma )
              \\
              =
              \overline{\sigma} \circ \pi_E( \tau_{n+1}( \tau_n( g ) ) ) \circ \overline{\sigma}
          \end{equation*}
          In conclusion, the following equation holds.
          \begin{align*}
              H( H( f ) )
              =
              \pi_E( \tau_{n+1}( \tau_n( g ) ) )
              =
              \overline{\sigma} \circ \pi_E( \tau_{n+1}( \tau_n( g ) ) ) \circ \overline{\sigma}
              =
              \overline{\sigma} \circ H( H( f ) ) \circ \overline{\sigma}
          \end{align*}
          Since $f$ was arbitrary, then $H( H( f ) ) = \overline{\sigma} \circ H( H( f ) ) \circ \overline{\sigma}$ for each $f \in ( P( \Sigma ) / E )( n, n )$.
    \item Let $f \in ( P( \Sigma ) / E )( n + 2, n + 2 )$ and $0 \le k \le n$.
          Then there exists a $g \in P( \Sigma )( n + 2, n + 2 )$ such that $\pi_E( g ) = f$.
          Then $H( f ) = H( \pi_E( g ) ) = F( g )$ the following equation holds.
          \begin{equation*}
              H( \gamma_{n,k} \circ f \circ \gamma_{n,k} )
              =
              H( \gamma_{n,k} \circ \pi_E( g ) \circ \gamma_{n,k} )
              =
              H( \pi_E( \gamma_{n,k} \circ g \circ \gamma_{n,k} ) )
              =
              F( \gamma_{n,k} \circ g \circ \gamma_{n,k} )
          \end{equation*}
          Since $\tau_{n+2}( \gamma_{n,k} \circ g \circ \gamma_{n,k} ) = ( 1 \boxtimes \gamma_{n,k} ) \circ \tau_{n+2}( g ) \circ ( 1 \boxtimes \gamma_{n,k} )$, then the following equation holds.
          \begin{align*}
              F( \gamma_{n,k} \circ g \circ \gamma_{n,k} )
              &=
              \pi_E( \tau_{n+2}( \gamma_{n,k} \circ g \circ \gamma_{n,k} ) )
              \\
              &=
              \pi_E( ( 1 \boxtimes \gamma_{n,k} ) \circ \tau_{n+2}( g ) \circ ( 1 \boxtimes \gamma_{n,k} ) )
              \\
              &= ( 1 \boxtimes \gamma_{n,k} ) \circ \pi_E( \tau_{n+2}( g ) ) \circ ( 1 \boxtimes \gamma_{n,k} )
              \\
              &= ( 1 \boxtimes \gamma_{n,k} ) \circ F( g ) \circ ( 1 \boxtimes \gamma_{n,k} )
              \\
              &= ( 1 \boxtimes \gamma_{n,k} ) \circ H( f ) \circ ( 1 \boxtimes \gamma_{n,k} )
          \end{align*}
          Since $f$ was an arbitrary morphism, then $H( \gamma_{n,k} \circ f \circ \gamma_{n,k} ) = ( 1 \boxtimes \gamma_{n,k} ) \circ H( f ) \circ ( 1 \boxtimes \gamma_{n,k} )$ for each $f \in ( P( \Sigma ) / E )( n + 2, n + 2 )$ and $0 \le k \le n$.
    \end{enumerate}
    Then $H$ is a control functor.
\end{proof}